\section{Discussão}\label{sec:discussao}

\paragraph{Validação como requisito.} O fato de o núcleo reproduzir $|G|$, as ordens de seis subgrupos, a distribuição completa de distâncias do \cube{2} e as primeiras esferas do \cube{3} não é um resultado novo, mas é o que dá credibilidade aos demais. Esses testes exercitam simultaneamente a geometria dos adesivos, a convenção de composição, a conversão para cubinhos, as tabelas de coordenadas e os algoritmos de busca: um erro de sinal numa rotação, por exemplo, alteraria a distribuição de distâncias. Derivar tudo de uma única definição geométrica (em vez de tabelas digitadas) reduziu o número de pontos em que tal erro poderia ocorrer.

\paragraph{Dificuldade não é comprimento de embaralhamento.} O E1 mostra que a distância real satura muito antes do comprimento do embaralhamento e se concentra em torno da média do grafo: no \cube{2}, 75\% dos estados estão a distância 8 ou 9 (\cref{fig:e0}), e embaralhamentos de 14 ou mais movimentos já têm $d^*$ médio praticamente igual à média uniforme. Experimentos que reportam desempenho ``por número de movimentos de embaralhamento'' medem, portanto, uma mistura de dificuldades. Para o \cube{2}, em que $d^*$ é conhecido, estratificar por $d^*$ e amostrar uniformemente cada esfera (E2, E3) elimina esse viés. No \cube{3}, em que $d^*$ é desconhecido em geral, a coincidência $d^*=L$ observada para $L\le12$ é uma propriedade das instâncias geradas (embaralhamentos curtos raramente admitem atalhos), não uma lei: ela não se sustenta para $L$ próximo do diâmetro 20~\cite{rokicki2014}.

\paragraph{Algoritmos.} A razão de crescimento da BFS (\gclEoneGrowthBFS) é próxima da razão entre esferas consecutivas do \cube{2} na região intermediária ($|S(6)|/|S(5)|\approx5{,}0$), e a busca bidirecional cresce aproximadamente com a sua raiz quadrada, como prevê a análise clássica~\cite{pohl1971,russell2021}. A$^*$ e IDA$^*$ ilustram o compromisso tempo--memória: IDA$^*$ expande mais nós (reexpansões entre iterações e ausência de detecção de repetidos), mas armazena apenas o caminho e, por ter custo por nó muito menor (sem tabela de dispersão nem heap), é o mais rápido neste problema --- o mesmo motivo pelo qual Korf~\cite{korf1997} o escolheu para o \cube{3}. Tempos absolutos, contudo, dependem da implementação e do hardware e devem ser lidos apenas comparativamente.

\paragraph{Heurísticas.} A eficiência acompanha a informatividade média $h(s_0)/d^*$, e o máximo de duas PDBs pequenas (729 e 5040 entradas) supera claramente cada uma delas, como esperado: as duas projeções capturam aspectos independentes do estado (orientação e permutação). O fator de ramificação efetivo de IDA$^*$ com $h_{\max}$ ($\approx2{,}9$ em $d^*=11$) contra $\approx3{,}4$--$3{,}8$ das PDBs isoladas resume esse ganho. A heurística de contagem de peças, embora correta, é quase inútil no \cube{2} (limitada a 2) --- um exemplo de que admissibilidade não implica utilidade. A heurística aditiva~\cite{felner2004} e PDBs maiores (por exemplo, a de todos os cantos do \cube{3}, com $8!\cdot3^7\approx8{,}8\cdot10^7$ entradas~\cite{korf1997}) são extensões naturais.

\paragraph{Ótimo versus válido.} No \cube{3}, o IDA$^*$ com nossas PDBs (cerca de 2 milhões de entradas) resolve otimamente instâncias de até 12 movimentos em segundos, mas o custo cresce cerca de uma ordem de grandeza a cada movimento adicional; estados aleatórios (tipicamente a 17--18 movimentos~\cite{rokicki2014}) estão fora de alcance com essa heurística. O método de duas fases, ao contrário, resolve qualquer estado em milissegundos com soluções em torno de 20 movimentos, mas sem garantia de mínimo --- distinção que a interface torna explícita. O fato de as soluções de duas fases terem encontrado o ótimo na maioria dos embaralhamentos curtos não altera essa natureza.

\paragraph{Estrutura e geradores.} O E4 ilustra como propriedades métricas de um grafo de Cayley dependem do conjunto gerador, e não apenas do grupo: o mesmo grupo de 29\,160 elementos tem diâmetro 17 ou 14 conforme se admitam meias-voltas, e geradores ``pobres'' como $\{\mv R,\mv{U2}\}$ produzem grafos longos (diâmetro 25 com 14\,400 vértices). As relações mínimas encontradas automaticamente ($X^4$, $XXX^2$, comutação de faces opostas, $(\mv{U2}\,\mv{R2})^6$) são exatamente as relações ``óbvias'' do cubo; relações mais longas e interessantes, como $(\mv{R U R' U'})^6=e$, aparecem como ciclos mais longos que o usuário pode fechar no grafo de estados.

\paragraph{Limitações.}
(i) A análise do GraphCube baseou-se no código-fonte; a aplicação publicada e a discussão no fórum não puderam ser acessadas no ambiente de desenvolvimento.
(ii) As buscas ótimas no \cube{3} usam PDBs modestas; não reproduzimos o desempenho de Korf nem os resultados de Rokicki et al.
(iii) A memória é medida pelo número de estados armazenados, não por bytes; no \cube{2} as tabelas de acesso direto têm tamanho fixo.
(iv) Os tempos foram medidos num único ambiente virtualizado e sem repetições independentes por instância; as comparações de tempo são indicativas.
(v) A suficiência da Primeira Lei foi verificada por amostragem, não demonstrada computacionalmente; a demonstração é a da literatura~\cite{joyner2008}.
(vi) A interface foi avaliada apenas pelos autores; não houve estudo com usuários.
